\section{Related work}
\label{sec:related}

\paragraph{Custom allocators and sanitizers.} CMASan~\cite{hong2025cmasan} is the work closest to
ours. It extends ASan to recognise a program's custom allocator functions and to give the objects they
return ASan-style protection, and it reports 19 previously missed bugs. CMASan
changes the tool. We keep the tools as shipped and measure how much their \emph{existing} annotation
interfaces and a rewrite recover, on bugs whose ground truth we control, and we include Memcheck and
HWASan, whose interfaces model memory differently from ASan's. Our results give a per-mechanism
account of the gap that CMASan closes: which misses come from missing per-object bounds (which
annotation fixes) and which come from the interface itself (which it does not).

\paragraph{Evaluating sanitizers.} Song et al.\ survey sanitizers and the guarantees each one gives
\cite{song2019sok}. MSET~\cite{vintila2025mset} evaluates the effectiveness of memory-safety
sanitizers systematically with a large set of test programs, and Jung et al.~\cite{jung2025practice}
compare memory-safety instrumentations for usability, performance and the guarantees they actually
provide. Test suites such as Juliet~\cite{nist2017juliet} and the cross-platform suite of Song et
al.~\cite{song2021testsuite} enumerate defect classes. These works vary the defect and the tool, but the
allocator the faulty object comes from is not a dimension of their evaluations. It is the one
dimension we vary. Our corpus follows the conventions of the bug corpora of the capability-hardware project
\emph{Capstone}~\cite{capstone} (one directory per defect, a negative control per bug,
positive controls that gate each run), from which we also took four real SQLite defects for a
preliminary transfer check (\S\ref{sec:tier2}).

\paragraph{Compiler effects.} Optimisations can remove the very access a sanitizer would have
checked when that access is undefined behaviour~\cite{isemann2023ub}. Our detection builds are
therefore unoptimised (\S\ref{sec:harness}), and only the performance measurements use \code{-O2}.

\paragraph{Allocators and hardware memory safety.} Memory tagging~\cite{serebryany2018mte} and
capability hardware such as CHERI~\cite{watson2015cheri} move bounds and lifetime checks into the
hardware, but the allocator still decides what an object is: a CHERI allocator must set capability
bounds and revoke stale capabilities itself~\cite{filardo2020cornucopia,bramley2023cheri}. Our
HWASan results show the same dependence in software: tags protect exactly the objects an allocator
chooses to tag, at the granularity the interface allows.
